\documentclass[tikz,border=6pt]{standalone}
% 공통 프리앰블 — PM Day1 교재 그림 (TikZ standalone)
\usepackage{fontspec}
\setmainfont{Apple SD Gothic Neo}
\setsansfont{Apple SD Gothic Neo}
\setmonofont{Menlo}
\usepackage{amssymb}
\usepackage{tikz}
\usetikzlibrary{arrows.meta,positioning,calc,shapes.geometric,fit,backgrounds,matrix,decorations.pathreplacing}
\definecolor{navy}{HTML}{1F3A5F}
\definecolor{teal}{HTML}{2A7F62}
\definecolor{rust}{HTML}{C8553D}
\definecolor{sand}{HTML}{E8E1D5}
\definecolor{ink}{HTML}{222222}
\definecolor{grey}{HTML}{7A7A7A}
\definecolor{lightgrey}{HTML}{EFEFEF}
\definecolor{navylight}{HTML}{DCE6F2}
\definecolor{teallight}{HTML}{DDF0E8}
\definecolor{rustlight}{HTML}{F6E0DA}
\tikzset{
  every node/.style={font=\small, text=ink},
  box/.style={draw=navy, thick, rounded corners=2pt, align=center, inner sep=5pt, fill=white},
  boxfill/.style={box, fill=navylight},
  boxteal/.style={draw=teal, thick, rounded corners=2pt, align=center, inner sep=5pt, fill=teallight},
  boxrust/.style={draw=rust, thick, rounded corners=2pt, align=center, inner sep=5pt, fill=rustlight},
  boxgrey/.style={draw=grey, thick, rounded corners=2pt, align=center, inner sep=5pt, fill=lightgrey},
  arr/.style={-{Stealth[length=2.5mm]}, thick, navy},
  arrteal/.style={-{Stealth[length=2.5mm]}, thick, teal},
  arrrust/.style={-{Stealth[length=2.5mm]}, thick, rust},
  arrgrey/.style={-{Stealth[length=2.5mm]}, thick, grey},
  lbl/.style={font=\footnotesize, text=grey, align=center},
  src/.style={font=\scriptsize, text=grey, align=left},
  title/.style={font=\normalsize\bfseries, text=navy, align=center},
}

\begin{document}
\begin{tikzpicture}[node distance=3mm]
\node[font=\small\bfseries, text=grey, anchor=south] at (22mm,3mm) {FROM};
\node[font=\small\bfseries, text=teal, anchor=south] at (86mm,3mm) {TOWARDS};
\node[font=\small\bfseries, text=navy, anchor=south] at (144mm,3mm) {Day1 교재의 대응 절};
\foreach \i/\n/\f/\t/\d in {
 1/{방향 1\\프로젝트 복잡성 이론}/{프로젝트와 PM의 라이프사이클 모델\\(단순한 순차 모델)}/{프로젝트와 PM의 복잡성에\\관한 이론들}/{1.4 NTCP 다이아몬드\\2.6 테일러링},
 2/{방향 2\\사회적 과정}/{도구적 라이프사이클 이미지 —\\프로젝트는 계획대로 실행되는 것}/{사회적 과정 — 사람들이\\상호작용하며 만들어 가는 것}/{1.3 임시조직론\\3.3 이해관계자},
 3/{방향 3\\가치 창출}/{제품 창출 — 산출물의 인도가 목표}/{가치 창출 — 편익과 가치가 목표}/{3.1 편익 관리 (ITO · BDN)},
 4/{방향 4\\프로젝트의 개념화}/{좁은 개념 — 사전에 정의된\\목표를 가진 것}/{넓은 개념 — 다중 목표, 다중\\이해관계자, 진행 중 정의가 바뀜}/{1.1 프로젝트 정의\\4.2 Morris},
 5/{방향 5\\실무자 발달}/{훈련받은 기술자\\trained technician}/{성찰적 실무자\\reflective practitioner}/{이 과정 전체가\\자격증 수업이 아닌 이유}}{
  \node[boxfill, text width=26mm, font=\footnotesize, anchor=north east, minimum height=14mm] (n\i) at (-4mm,{-(\i-1)*17mm}) {\n};
  \node[boxgrey, text width=44mm, font=\footnotesize, anchor=north, minimum height=14mm] (f\i) at (22mm,{-(\i-1)*17mm}) {\f};
  \node[boxteal, text width=44mm, font=\footnotesize, anchor=north, minimum height=14mm] (t\i) at (86mm,{-(\i-1)*17mm}) {\t};
  \node[box, draw=navy, fill=white, text width=44mm, font=\footnotesize, anchor=north, minimum height=14mm] (d\i) at (144mm,{-(\i-1)*17mm}) {\d};
  \draw[arrteal, thick] (f\i) -- (t\i); \draw[arrgrey, dashed] (t\i) -- (d\i);
}
\node[box, draw=rust, fill=rustlight, text width=160mm, font=\footnotesize, anchor=north west, align=left] at (-30mm,-88mm) {\textbf{흔한 오독} — 다섯 방향은 \textbf{병렬적} 방향이며 순차적 성숙 단계가 아니다. RPM은 방법론 교체("애자일이 waterfall을 대체")가 아니라 연구 대상 자체(라이프사이클 모델 → 임시조직 이론)의 교체다. PMBOK 7판(2021)의 원칙 전환은 이 논문(2006)의 15년 늦은 반영이고, 8판(2025)의 부분 회귀는 학술과 실무의 재타협이다.};
\node[src, anchor=north west] at (-30mm,-106mm) {출처: Winter, Smith, Morris \& Cicmil (2006), Directions for future research in project management, IJPM 24(8), Table 1을 바탕으로 재구성.};
\end{tikzpicture}
\end{document}
